% Options for packages loaded elsewhere
% Options for packages loaded elsewhere
\PassOptionsToPackage{unicode}{hyperref}
\PassOptionsToPackage{hyphens}{url}
\PassOptionsToPackage{dvipsnames,svgnames,x11names}{xcolor}
%
\documentclass[
  numafflabel,
  authoryear,
  preprint]{elsarticle}
\usepackage{xcolor}
\usepackage{amsmath,amssymb}
\setcounter{secnumdepth}{-\maxdimen} % remove section numbering
\usepackage{iftex}
\ifPDFTeX
  \usepackage[T1]{fontenc}
  \usepackage[utf8]{inputenc}
  \usepackage{textcomp} % provide euro and other symbols
\else % if luatex or xetex
  \usepackage{unicode-math} % this also loads fontspec
  \defaultfontfeatures{Scale=MatchLowercase}
  \defaultfontfeatures[\rmfamily]{Ligatures=TeX,Scale=1}
\fi
\usepackage{lmodern}
\ifPDFTeX\else
  % xetex/luatex font selection
\fi
% Use upquote if available, for straight quotes in verbatim environments
\IfFileExists{upquote.sty}{\usepackage{upquote}}{}
\IfFileExists{microtype.sty}{% use microtype if available
  \usepackage[]{microtype}
  \UseMicrotypeSet[protrusion]{basicmath} % disable protrusion for tt fonts
}{}
\makeatletter
\@ifundefined{KOMAClassName}{% if non-KOMA class
  \IfFileExists{parskip.sty}{%
    \usepackage{parskip}
  }{% else
    \setlength{\parindent}{0pt}
    \setlength{\parskip}{6pt plus 2pt minus 1pt}}
}{% if KOMA class
  \KOMAoptions{parskip=half}}
\makeatother
% Make \paragraph and \subparagraph free-standing
\makeatletter
\ifx\paragraph\undefined\else
  \let\oldparagraph\paragraph
  \renewcommand{\paragraph}{
    \@ifstar
      \xxxParagraphStar
      \xxxParagraphNoStar
  }
  \newcommand{\xxxParagraphStar}[1]{\oldparagraph*{#1}\mbox{}}
  \newcommand{\xxxParagraphNoStar}[1]{\oldparagraph{#1}\mbox{}}
\fi
\ifx\subparagraph\undefined\else
  \let\oldsubparagraph\subparagraph
  \renewcommand{\subparagraph}{
    \@ifstar
      \xxxSubParagraphStar
      \xxxSubParagraphNoStar
  }
  \newcommand{\xxxSubParagraphStar}[1]{\oldsubparagraph*{#1}\mbox{}}
  \newcommand{\xxxSubParagraphNoStar}[1]{\oldsubparagraph{#1}\mbox{}}
\fi
\makeatother


\usepackage{longtable,booktabs,array}
\usepackage{calc} % for calculating minipage widths
% Correct order of tables after \paragraph or \subparagraph
\usepackage{etoolbox}
\makeatletter
\patchcmd\longtable{\par}{\if@noskipsec\mbox{}\fi\par}{}{}
\makeatother
% Allow footnotes in longtable head/foot
\IfFileExists{footnotehyper.sty}{\usepackage{footnotehyper}}{\usepackage{footnote}}
\makesavenoteenv{longtable}
\usepackage{graphicx}
\makeatletter
\newsavebox\pandoc@box
\newcommand*\pandocbounded[1]{% scales image to fit in text height/width
  \sbox\pandoc@box{#1}%
  \Gscale@div\@tempa{\textheight}{\dimexpr\ht\pandoc@box+\dp\pandoc@box\relax}%
  \Gscale@div\@tempb{\linewidth}{\wd\pandoc@box}%
  \ifdim\@tempb\p@<\@tempa\p@\let\@tempa\@tempb\fi% select the smaller of both
  \ifdim\@tempa\p@<\p@\scalebox{\@tempa}{\usebox\pandoc@box}%
  \else\usebox{\pandoc@box}%
  \fi%
}
% Set default figure placement to htbp
\def\fps@figure{htbp}
\makeatother





\setlength{\emergencystretch}{3em} % prevent overfull lines

\providecommand{\tightlist}{%
  \setlength{\itemsep}{0pt}\setlength{\parskip}{0pt}}



 
\usepackage[]{natbib}
\bibliographystyle{elsarticle-harv}


\makeatletter
\@ifpackageloaded{caption}{}{\usepackage{caption}}
\AtBeginDocument{%
\ifdefined\contentsname
  \renewcommand*\contentsname{Table of contents}
\else
  \newcommand\contentsname{Table of contents}
\fi
\ifdefined\listfigurename
  \renewcommand*\listfigurename{List of Figures}
\else
  \newcommand\listfigurename{List of Figures}
\fi
\ifdefined\listtablename
  \renewcommand*\listtablename{List of Tables}
\else
  \newcommand\listtablename{List of Tables}
\fi
\ifdefined\figurename
  \renewcommand*\figurename{Figure}
\else
  \newcommand\figurename{Figure}
\fi
\ifdefined\tablename
  \renewcommand*\tablename{Table}
\else
  \newcommand\tablename{Table}
\fi
}
\@ifpackageloaded{float}{}{\usepackage{float}}
\floatstyle{ruled}
\@ifundefined{c@chapter}{\newfloat{codelisting}{h}{lop}}{\newfloat{codelisting}{h}{lop}[chapter]}
\floatname{codelisting}{Listing}
\newcommand*\listoflistings{\listof{codelisting}{List of Listings}}
\makeatother
\makeatletter
\makeatother
\makeatletter
\@ifpackageloaded{caption}{}{\usepackage{caption}}
\@ifpackageloaded{subcaption}{}{\usepackage{subcaption}}
\makeatother
\journal{Journal Name}
\usepackage{bookmark}
\IfFileExists{xurl.sty}{\usepackage{xurl}}{} % add URL line breaks if available
\urlstyle{same}
\hypersetup{
  pdftitle={Predictive models of fire severity: from complex distribution modelling to machine learning},
  pdfauthor={Ceres Barros; Cameron Naficy; Steve Cumming; David W. Andison; Eliot J. B. McIntire},
  colorlinks=true,
  linkcolor={blue},
  filecolor={Maroon},
  citecolor={Blue},
  urlcolor={Blue},
  pdfcreator={LaTeX via pandoc}}


\setlength{\parindent}{6pt}
\begin{document}

\begin{frontmatter}
\title{Predictive models of fire severity: from complex distribution
modelling to machine learning}
\author[1,2]{Ceres Barros%
\corref{cor1}%
}

\author[1,3]{Cameron Naficy%
%
}

\author[4]{Steve Cumming%
%
}

\author[5]{David W. Andison%
%
}

\author[1,2,4]{Eliot J. B. McIntire%
%
}


\affiliation[1]{organization={Canadian Forest Service (Pacific Forestry
Centre), Natural Resources Canada},addressline={506 Burnside Road
West},city={Victoria},country={Canada},countrysep={,},postcode={V8Z
1M5},postcodesep={}}
\affiliation[2]{organization={Department of Forest Resources Management,
University of British Columbia},addressline={2424 Main
Mall},city={Vancouver},country={Canada},countrysep={,},postcode={V6T
1Z4},postcodesep={}}
\affiliation[3]{organization={USDA Forest Service, Pacific Northwest
Region},city={Baker City},country={USA},countrysep={,},postcodesep={}}
\affiliation[4]{organization={Faculté de Foresterie, de Géographie et de
Géomatique, Département des Sciences du Bois et de la Forêt, Pavillon
Abitibi-Price, Université Laval},addressline={2405, rue de la
Terrasse},city={Québec},country={Canada},countrysep={,},postcode={G1V
0A6},postcodesep={}}
\affiliation[5]{organization={fRI Research},addressline={1176 Switzer
Drive},city={Hinton},country={Canada},countrysep={,},postcode={T7V
1V3},postcodesep={}}

\cortext[cor1]{Corresponding author}





        
\begin{abstract}
Wildfire is a common disturbance in many forest systems across the
globe, where it influences a variety of ecosystem processes and
functioning, as well as community safety. In particular, fire severity
(i.e.~fire-caused vegetation mortality) affects post-fire vegetation
succession, ecosystem recovery pathways and future fire regimes. With
wildfires becoming more frequent and severe, our ability to predict fire
severity is increasingly necessary to properly forecast changes in fire
regimes, forest structure and resilience, and to allocate adequate
wildfire and forest management strategies to protect ecosystems and
communities. Fire severity, however, can be difficult to model and
predict. Apart from data availability issues, fire severity responds to
multiple drivers (climate, topography and fuel composition covariates)
and can present statistical hurdles. Amongst these is inflation at the
tails when severity is measured as a proportion variable (i.e.~\%
mortality). We explore the performance of two statistical approaches to
model and predict fire severity: a frequentist approach using a
zero-and-one-inflated beta (ZOIB) regression fitted within a GAMLSS
framework, and a machine learning approach using extreme gradient
tree-boosting using XGBoost. While the ZOIB regression allowed modelling
the influence of different biophysical covariates on distinct aspects of
severity -- i.e the probability of no severity (0\% mortality), the
probability of stand-replacement (100\% mortality), and the average and
variance of severity -- it came at the cost of speed and a long process
of model selection. On the other hand, XGBoost showed high speed
performance and slightly higher fit, at the cost of interpretability and
the inability to deal with overdispersion at the tails and with random
effects. Interestingly, after accounting for the random effect of fire
ID, the ZOIB regression did not support including weather-related
variables in the model, which were given relatively high importance in
XGBoost models. Maps of predicted fire severity revealed high agreement
between the ZOIB regression and XGBoost and good ability to predict
intermediate fire severity values. We conclude that while ZOIB
regression frameworks allow modelling severity within a given fire
(i.e.~irrespective of between-fire differences) and a finer
interpretation of its drivers, XGBoost keeps high accuracy and provides
the ability to swiftly update predictions using large datasets, which is
of importance in a wildfire management context given yearly updates to
fire data streams.
\end{abstract}





\end{frontmatter}
    

\section{Introduction}\label{introduction}

Wildfires are an integral part of several ecosystems worldwide {[}CITE:
Bowman-2009{]}, including Canadian montane and boreal forests {[}REF{]},
but can cause important biodiversity and socio-economic disruptions.
Wildfires affect numerous ecosystem processes such as ecosystem recovery
pathways {[}REF{]}, permafrost degradation {[}REF{]}, carbon emissions,
as well as biodiversity more generally {[}REF{]}, community health and
safety and infrastructure {[}REF{]}. The dimension of these impacts
depends on ecological and socio-economic context, but also on fire
regime properties, like size, frequency and severity. Given ongoing
climate warming trends in Canada and elsewhere, we expect changes in
future fire regimes to include increases in all these three dimensions.
Predicting and quantifying these changes will be paramount to managing
fire-disturbed forests under future climate scenarios.

Modelling and forecasting fire severity -- here defined as fire-caused
aboveground vegetation mortality -- can be particularly challenging.
Apart from data scarcity in drivers of severity and severity itself,
particularly for fires pre-dating remote sensing, modelling severity can
present several statistical hurdles. The inherent variability of fire
behaviour and the nature of severity data both decrease predictive power
at non-extreme severity values, and have challenged previous attempts at
classifying fire perimeters into severity classes. First, fire severity
is the result of a complex interplay between topography, climate and
fuel conditions, exhibiting high spatial variability between and within
fire events {[}CITE: Andison-2014; CITE: Viedma-2020{]} and non-linear
responses to these drivers. For instance, extreme climate or weather
conditions are generally expected to override the effects of fuel
composition {[}CITE: Cruz-2022; CITE: Miquelajauregui-2016{]}. Different
drivers also affect different fire characteristics distinctively and at
different spatio-temporal scales {[}REF{]}, with weather exerting a
stronger effect on fire size than fuel composition, which likely has a
stronger impact on finer scale severity patch configuration {[}REF{]}.

Second, severity measurements can have some degree of uncertainty,
especially when based on aerial photo interpretation and remotely-sensed
variables {[}CITE: Keeley-2009{]}, and can be measured using different
approaches for different fire events. To circumvent these issues and
facilitate comparative analyses of severity, severity is often reported
in classes {[}CITE: Miller-2007{]}. However, using categorical or
ordinal variables can lead to loss of information and predictive power
{[}CITE: Kuhn-2013{]} and introduce estimation bias. Furthermore, when
modelling severity classes (e.g.~severity classes 0 (0\% mortality), 1
(1-25\%), 2 (26-50\%), 3 (51-75\%), 4 (76-94\%) and 5 (\$\geq\$95\%
mortality)) using a categorical/nominal variable instead of an ordinal
scale leads to i) the violation of the assumption of equality of
intervals between classes (i.e.~the distance between classes 0 and 1 is
smaller than between classes 0 and 4), and consequently of the normality
of residuals, and to ii) a mismatch between the scales of the latent
variable (i.e.~the ``true'', underlying, severity ranging from 0-100\%
mortality) and the measured variable (i.e.~severity in classes 0-5).
Both issues can cause bias and error in parameter estimation, and
failure in rejecting the null hypothesis (see {[}CITE: Harwell-2001{]}
for a review of these issues). Using a continuous variable of severity
should therefore improve models, provided that the scale of measurement
is consistent between fire events and across ecosystems {[}REF{]}.

Continuous severity metrics are often measured in relative terms
(i.e.~proportion of burnt or dead vegetation, normalized burn ratios and
relativized CBI; {[}REF{]}), which brings other statistical
considerations. Proportion data is bounded in the 0-1 interval and in
the case of severity both extremes (zeroes and ones) are possible values
with a biophysical meaning and should therefore not be excluded or
transformed. Zero-one inflated beta distributions are suitable to
describe proportion data in the closed interval {[}0,1{]} {[}CITE:
Ospina-2010{]}. It consists of a mixture of a continuous beta
distribution on the (0,1) interval and a distribution assigning a
non-negative probability to zeros and/or ones, like the Bernoulli
distribution, allowing to model severity data in the {[}0,1{]} range.
Zero-and-one-inflated beta regressions are described by four parameters:
the mean (\(\mu\)) and the precision (\(\sigma\), inverse of dispersion)
parameters of the continuous component (related to the shape and scale
parameters of the beta distribution) and two parameters, \(\nu\) and
\(\tau\), linked to the probability of observing a zero (p0) and the
probability of observing a one (p1), respectively. All four parameters
can be modelled with unique sets of predictors in a regression framework
{[}CITE: Masserini-2017; CITE: Smithson-2013{]}. Generalised additive
models of location, scale and shape (GAMLSS) provide an incredibly
flexible tool to implement zero--and-one-inflated beta (ZOIB)
regressions using a GLM-type approach {[}CITE: Stasinopoulos-2017{]} and
explore how different aspects of fire severity are affected by
environmental and vegetation drivers.

Frequentist approaches like the GAMLSS are, however, often surpassed by
machine learning in terms of predictive accuracy and their ability to
ingest large amounts of data {[}REF{]}. As wildfire-related data grows,
especially data derived from remote sensed variables, forecasting models
need to be able to rapidly update predictions with ever increasing data
volumes. Extreme gradient boosting (XGBoost; {[}CITE: Chen-2016{]}) is a
highly performant machine learning technique, used across several
fields, including medicine {[}CITE: Diez-Sanmartin-2021{]}, economy
{[}CITE: Chalvatzis-2020{]}, computer sciences {[}CITE: Cui-2021{]},
ecology {[}CITE: Batunacun-2021{]} and others. It often provides higher
accuracy than other methods {[}REF{]}, can deal with complex non-linear
interactions between predictors (`features'), various data classes
{[}REF{]} and both classification and regression problems {[}CITE:
Wiens-2025{]}.

Here, we use ZOIB models in a GAMLSS framework to develop a predictive
model of fire severity that overcomes the statistical issues associated
with proportional fire severity data. Our goal is to provide a better
understanding of how fire severity responds to pre-fire vegetation
(i.e.~fuel), weather and topographic covariates, while taking into
account its continuous nature to ultimately improve landscape-level
predictions of fire severity. For this, we used an existing high
resolution photo-based historical wildfire dataset that covers a large
spatial extent across western boreal Canada, and where severity was
measured as the \% of burnt/dead vegetation following fire using a
consistent approach across fire events. These data were paired with
pre-fire vegetation, land-cover, fire weather and topography data
{[}CITE: Ferster-2016{]}. We explored how average severity, no severity
and stand-replacement (0\% and 100\% mortality, or the probabilities of
zeroes and ones in ZOIB), as well as the variance of severity may be
affected by distinct covariates. Drawing from previous evidence that
extreme fire weather can override fuel and topography effects on driving
high severity fires and that the lack of fuel can act as a fire barrier,
we hypothesised that 1) weather variables would show a stronger
influence on the probability of stand-replacement than topography,
vegetation and land-cover variables, while 2) vegetation, land-cover and
topography covariates would have a stronger influence on average
severity values and their variability, and 3) vegetation and land-cover
would be the strongest drivers of minimum values of severity (0\%
mortality).

\section{Materials and Methods}\label{materials-and-methods}

We developed a ZOIB regression model of fire severity using the fire
severity, vegetation, land-cover, topography and weather data described
by {[}CITE: Ferster-2016{]}. Although the original datasets were
identical, we processed several variables differently and the final
suite of candidate variables differed substantially from Ferster et
al.'s {[}CITE: Ferster-2016{]}. Here, we provide a summary of these data
and how they were processed for the present work.

\subsection{Study area}\label{study-area}

The data spans a broad study area covering sections of the Boreal Plain,
Taiga Plain, Boreal Shield and Montane Cordillera ({[}REF{]}; Ecozones
Canada), in the Alberta and Saskatchewan provinces. Data in the Boreal
and Taiga Plain are on the Boreal Forest and Foothills regions
({[}REF{]}; NR Alberta); data in the Boreal and Taiga Shield are on the
Canadian Shield region; and data in the Montane Cordillera are on the
Rocky Mountain region (Figure 1). Hence, the data covers a large
gradient of vegetation, weather and topography conditions {[}CITE:
Ferster-2016{]}. The Boreal Plains (BP) are predominantly flat or with
gently rolling hills, where north to south (mixed boreal forest
transitioning to broadleaf dominated stands) and uplands to lowlands
(transitions from dry to wet conditions and associated species)
vegetation composition gradients exist concurrently with true wetlands.
The landscape in the Canadian Shield (CS) is more rugged, with thinner
soils and also abundant wetlands. Pure conifer and mixed stands dominate
the forested areas. The Foothills (Fh) differs from the BP thanks to its
steeper slopes and higher elevations, which create better conditions for
lodgepole pine. The Rocky Mountain (RM) region is dominated by montane
landscapes, with strong slope and elevation gradients, which generate a
variety of local conditions that lead to distinct vegetation types at
relatively small scales. A more detailed description of each region's
vegetation composition is available in {[}CITE: Ferster-2016{]}.

\subsection{Fire severity data}\label{fire-severity-data}

Fire severity data was obtained from a historical wildfire photo-based
dataset initially composed of 129 fires, of which 36 had available
pre-fire vegetation data and met the following criteria: i) no evidence
of pre-fire anthropogenic activity, ii) minimal or no fire suppression,
no post-fire salvage logging and iii) high quality, high resolution
aerial photos available within five years post-fire (1:20,000 or large
format plates greater than 1:31,860 of sufficient quality to resolve
individual trees). All photos were interpreted by the same person and
fire perimeters were subdivided into six classes of post-fire vegetation
mortality at 10m² resolution: class 0: no mortality, 0\%; class 1: 1\%
to 25\%; class 2: 26\% to 50\%; class 3: 51\% to 75\%; class 4: 76\% to
94\%; and class 5: \$\geq\$95\%. Vegetation mortality included percent
dead tree crowns in forested areas, dead shrubs in non-forested areas
with woody vegetation, and percent area scorched in non-forested areas
of grass or bryophytes/lichen. The resulting digital mortality maps were
saved at 1:50,000 resolution. Please see {[}CITE: Andison-2012{]} and
{[}CITE: Ferster-2016{]} for further detail.

Here, we converted severity to a continuous variable as an attempt to i)
improve predictive power and ii) to overcome aforementioned residual
assumption and bias issues related to ordinal variables. We used the
mean percent value of each class as the new, continuous, severity value.
We also buffered fire perimeters inwards and outwards to include fire
islands (following {[}CITE: Andison-2014{]}) -- 0\% mortality -- to
better explore biotic and abiotic conditions leading to
no-fire-severity.

\subsection{Pre-fire vegetation data}\label{pre-fire-vegetation-data}

Pre-fire vegetation data was obtained from high resolution photos
obtained within five years pre-fire and interpreted using the Alberta
and Saskatchewan vegetation inventory methods. Vegetation was
interpreted to homogeneous polygons and covered all locations within a
100m outer buffer of the fire perimeters. Interpreted vegetation
variables included: percent crown cover, species composition, height
class, and age class for forested areas, vegetation type for
non-forested areas, and soil moisture class. Height class, age class and
crown cover class were converted to continuous variables using the mean
value of each class. Since the vegetation data structure differed
depending on the standard used (Alberta vs Saskatchewan), we harmonised
it to the Common Attribute Schema for Forest Resource Inventories
(CASFRI) standard {[}CITE: Cosco-2011{]}. Unlike {[}CITE:
Ferster-2016{]}, we modelled fire severity against continuous vegetation
composition covariates (see Table 1), instead of fuel types derived from
vegetation composition. Only overstory tree species cover variables
(`layer' 1 of the CASFRI standard) were considered, as understorey
species identification and quantification from aerial photography is
prone to error. Overstory species present across the data were: white
(\emph{Picea glauca} (Moench) Voss), black (\emph{P. mariana} (Mill.)
B.S.P.) and Engelmann spruce (\emph{P. engelmannii} Parry ex Engelm.),
lodgepole pine (\emph{Pinus contorta} (Douglas) var. \emph{latifolia}
(Engel.) Critchfield), jack pine (\emph{P. banksiana} Lamb.), tamarack
(\emph{Larix laricina} (Du Roi) K. Koch), balsam fir (\emph{Abies
balsamea} (L.) Rich.), subalpine fir (\emph{A. lasiocarpa} (Hooker)
Nuttall), trembling aspen (\emph{Populus tremuloides} Michx.), balsam
poplar (\emph{Populus balsamifera} L.) and paper birch (\emph{Betula
papyrifera} Marshall). Understorey species \% cover was collapsed into a
single binary variable (`understorey'), indicating presence/absence of
any species in the understorey (i.e.~understorey cover \textgreater{}
0\%). A binary variable, `isForest', was created to distinguish between
forest polygons and non-forest polygons (both vegetated and
non-vegetated). ``Forested'' and ``wooded'' wetlands were treated as
forested polygons. Preliminary analyses of fire severity responses to
covariates (see below for detail) suggested that soil moisture regime
classes and species could be grouped to avoid covariate redundancy.
Hence, we created two soil moisture regime (SMR) variables based on the
original four-class categorical CASFRI variable -- a categorical ordered
variable of classes 1 to 5 (from driest to wettest) and a binary
variable (SMRwet) coding the two wettest classes as `1' and as other
classes as `0' (Table 1). Species were grouped into flammable conifers
(`FlamConif') that included spruces (\emph{Picea} spp.) and pines
(\emph{Pinus} spp.), an \emph{Abies} spp. group (`Abies'), and a
broadleaf species group (`Broadleaf') including aspen, poplar species
(\emph{Populus} spp.) and paper birch (\emph{Betula papyrifera}).
Species group' percent cover values were calculated as the sum of their
species percent cover values.

\subsection{Fire weather data}\label{fire-weather-data}

We used the historic fire weather data described in {[}CITE:
Ferster-2016{]} and {[}CITE: Amiro-2001{]}, which consisted of daily
temperature, relative humidity, wind speed and rain precipitation
averaged from weather station measurements during the first 21 days
following fire ignition date. These data were used to calculate six fire
weather indices (using the Canadian Forest Fire Weather Index and Fire
Behaviour Prediction System; {[}CITE:
Forestry-Canada-Fire-Danger-Group-1992{]}): fine fuel moisture code
(FFMC), duff moisture code (DMC), drought code (DC), buildup index
(BUI), initial spread index (ISI) and fire weather index (FWI). Weather
variables and fire weather indices were then summarised, per fire, using
the minimum, maximum, mean, range, and coefficient of variation across
the whole period of fire weather data (up to 21 days) and across the
first three days of fire weather, or using only the first day of fire
weather.

\subsection{Topography and other data}\label{topography-and-other-data}

Six topographic variables were calculated from a digital elevation map
at 30m resolution: elevation, percent slope, aspect, two terrain
roughness variables (surface relief ratio, `SRR', and surface aspect
ratio, `SAR') and a soil wetness variable (compound topographic index,
`CTI'; see {[}CITE: Ferster-2016{]} for further detail). In addition, we
included variables coding Canadian ecoregions and Alberta ecoregions
(see Table 1). To account for the neighbourhood (or contagion) effects
of fire severity, we generated a proxy variable of neighbourhood
influence (`ngbPropBurns') calculated as the proportion of burnt
neighbours of each focal cell in buffers of different sizes (30m, 60m,
120m, 240m, 480m, 960m, 1920m radiuses). We chose to calculate the
proportion of burn neighbours instead of using neighbour severity
values, because we intend this model to be used for predicting fire
severity and so explanatory variables cannot depend on knowing values of
severity (the response variable) a priori. Finally, we included fire
name (`FireID') as a categorical variable.

Fire severity and vegetation polygons were gridded (i.e.~rasterized)
using a 30m cell size, matching the topographic data layers, and all
three datasets were then converted to point shapefile data to speed up
their spatial intersection. Because the original datasets were projected
using different coordinate reference systems (CRS), the final dataset
had minor alignment artefacts. These are unlikely to affect our
analyses. Also, points where one of the data layers had missing
information were excluded from further analyses. Fire weather data was
i.e.~a-spatial and homogenous within a fire, as each fire had only a
single value per fire weather variable.

\subsection{Zero-and-one-inflated beta regression
models}\label{zero-and-one-inflated-beta-regression-models}

Beta regressions are similar to generalised linear models (GLMs;
{[}CITE: Bolker-2009{]}; {[}CITE: McCullagh-Nelder-1989{]}) in that they
are composed of three components: the random component (the beta
distribution), the systematic component (the linear predictor) and the
link function (the link between the random and systematic component;
please see {[}CITE: Ferrari-2004{]} for a comprehensive explanation).
The beta distribution is a member of the exponential family defined by
two parameters, a and b (both bound to (0,1)), that allow for a variety
of shapes (unimodal, `U', `J', inverted `J', uniform; {[}CITE:
Ospina-2010{]}). In a regression context, the \(\mu\) (expected mean)
and \(\sigma\) parameterisation, Beta(\(\mu\), \(\sigma\)), is used
instead of Beta(a,b) (see {[}CITE: Douma-2019{]} for how the two
parameterisations are related). Using the Beta(\(\mu\), \(\sigma\))
parameterisation, the expected mean \(\mu\) (E(Y) = \(\mu\)) is related
to the expected variance via the precision parameter \(\sigma\) {[}CITE:
Ospina-2010{]}:

\begin{equation}\phantomsection\label{eq-1}{
\text{var}(y) = \frac{\mu (1 - \mu)}{1 + \sigma}
}\end{equation}

The most commonly used link function for the beta regression is the
logit function (Equation~\ref{eq-2} and Equation~\ref{eq-3}), which
ensures that values from the linear predictor (\(\eta\) in
Equation~\ref{eq-3}) fall between (0,1). Best predictor coefficients
(\(X\beta\)s in Equation~\ref{eq-3}) and associated standard errors and
confidence intervals are obtained by maximum likelihood estimation.

\begin{equation}\phantomsection\label{eq-2}{
y \sim \text{Beta}(\mu, \sigma)
}\end{equation}

\begin{equation}\phantomsection\label{eq-3}{
\text{logit}(\mu) = \log\left(\frac{\mu}{1 - \mu}\right) = \eta = X\beta
}\end{equation}

The \(\sigma\) parameter can also be modelled with covariates using a
log link function. We refer the reader to {[}CITE: Douma-2019{]} for an
accessible and detailed explanation of the beta regression.

Because the beta distribution is defined on the interval (0,1),
modelling proportion data that includes zeros and ones requires
expanding the beta distribution and using a zero-and-one-inflated beta
(ZOIB) regression {[}CITE: Ospina-2010{]}. The ZOIB regression assumes
that y arises from two Bernoulli processes describing the probability of
observing a zero and the probability of observing a one, and a beta
regression model describing all non-zero and non-one observations. The
ZOIB is then defined as a four-parameter distribution, ZOIB(\(\mu\),
\(\sigma\), \(\nu\), \(\tau\)), where \(\nu\) and \(\tau\) are related
with the probabilities of observing a zero (P(Y=0), p0) and one (P(Y=1),
p1):

\begin{equation}\phantomsection\label{eq-4}{
p0 = \frac{\nu}{1 + \nu + \tau}
}\end{equation}

\begin{equation}\phantomsection\label{eq-5}{
p1 = \frac{\tau}{1 + \nu + \tau}
}\end{equation}

It is worth noting that when using a ZOIB regression \(\mu\) and
Var(\(\mu\)) are no longer the expected mean (E(Y)) and variance of the
complete dataset. The calculation of the expected rth value, E(Yr), from
a ZOIB is a weighted average of the rth moment (i.e.~realisation) of the
Bernoulli distributions and the corresponding moment of the beta
distribution {[}CITE: Ospina-2010; CITE: Stasinopoulos-2017{]}:

\begin{equation}\phantomsection\label{eq-6}{
E(Y^r) = \frac{\hat{\mu} + \hat{\tau}}{1 + \hat{\nu} + \hat{\tau}}
}\end{equation}

where \(\hat{\mu}\), \(\hat{\nu}\) and \(\hat{\tau}\) are fitted
parameter values {[}CITE: Stasinopoulos-2007{]}. For the variance of the
ZOIB please see Appendix 1.

Using the GAMLSS framework implemented in the \texttt{gamlss} R package
{[}CITE: Stasinopoulos-2007{]}, we fit a ZOIB regression model to
analyse the response of fire severity in terms of its mean, variance,
probability of observing no severity (zeros, 0\% mortality) and
probability of observing maximum severity (ones, 100\% mortality). We
modelled each parameter of the ZOIB (\(\mu\), \(\sigma\), \(\nu\),
\(\tau\)) as a function of different sets of covariates {[}CITE:
Stasinopoulos-2017{]}.

Given the large number of potential drivers of fire severity (Table 1)
we ran preliminary multivariate and correlation analyses to filter
candidate variables (Appendix 2). Principal components analyses (PCAs)
on vegetation, topography and weather covariates, on which we overlaid
fire severity (using the \texttt{env.fit} function in the \texttt{vegan}
R package; {[}REF{]}) revealed a very noisy dataset and no visually
obvious relationships between covariates and severity. However, they
showed groupings by FireID and SMR. Spearman Rank correlations showed
relatively low correlation values (rho \textless{} \textbar0.5\textbar)
between most vegetation covariate pairs -- note that given the size of
the data (n=437,557) interpreting p-values is meaningless as all
correlation tests are significant at p\textless0.05. Exceptions were
tamarack and height (rho = -0.53), tamarack and SMR (0.58), SMR and
crown closure (-0.54), and SMR and height (-0.73). Correlations were
also low between vegetation and topography covariates (the exception
being elevation and jack pine, with rho = -0.63) and between topography
covariates (except CTI and slope, rho = -0.68, and CTI and SAR, -0.68).
Conversely, weather covariates were often highly correlated (rho
\textgreater{} 0.8). Hence, we selected the vegetation, topography and
weather covariates with highest correlations with fire severity, but
avoided choosing collinear weather covariates (see bold variables in
Table 1). Despite their low correlations with severity, the vegetation
variables Abies, Broadleaf and height were kept as potential candidates,
while crown closure and origin (i.e.~year of origin of the stand) were
removed due to lack of data for some fires. We also kept both SMR and
SMRwet, all topographic variables and all ngbPropBurn variables
(differing in buffer size). As for weather variables, highest
correlations with severity were obtained for fire weather indices and
wind speed, but the latter was dropped since it enters in the
calculation of fire weather indices. For each fire weather index we kept
the mean and coefficient of variation, as the latter should be able to
summarise the variability of weather conditions given by the minimum,
maximum and range, and fitted separate models using 1, 1-3 or 21 days of
fire weather.

This preliminary selection of variables resulted in eight potential
vegetation covariates, six topography covariates, 15 fire weather
covariates (for any given number fire weather days) and seven
ngbProbBurn covariates. We proceeded to fit ZOIB regressions using a
stepwise approach based on Generalised Akaike Information Criterion
(GAIC) scores, penalised by the log-number of observations. Following
{[}CITE: Stasinopoulos-2017{]}, we began by doing forward and backward
stepwise selection on the \(\mu\) model (i.e., non-zeros and non-ones),
starting from a null model and capping the upper scope of the selection
to include all main effects. After finding the \(\mu\) model with the
lowest GAIC, we repeated this process for the \(\sigma\) model (using
the best \(\mu\) model), then the \(\nu\) model and finally the \(\tau\)
model. We then included pairwise interaction terms between covariates of
the same type (except species cover covariates, as we expected their
effects to be mostly additive). Including interactions between different
types of variables or more than two variables led to convergence issues
and overly complex models that became hard to interpret. Given that our
goal was to develop a predictive model of fire severity independent of
conditions specific to any particular fire, or location, beyond those
already explained by the measured biophysic covariates, we included Fire
ID as a random effect on the intercepts of all parameter models. In this
first round of model selection, we attempted to include fire weather
variables summarised across the first 21 days of weather. No fire
weather indices had significant effects on fire severity, so we repeated
the process using weather variables summarised across days 1-3 and using
only the first day of weather data. In both cases, fire weather indices
did not have significant effects or improve the models and were excluded
from further analyses.

We proceeded to simplify the models by excluding covariates and, or,
covariate interactions with non-significant effects. Although this
reduced model fit in terms of GAIC score, we valued model parsimony and
interpretability above optimising GAIC scores. After selecting the final
biophysical covariates in each model, we ran another forward stepwise
selection process to include neighbourhood covariates at different
scales (ngbProbBurn variables, Table 1). Again, this was done
sequentially for the \(\mu\), \(\sigma\), \(\nu\) and \(\tau\) models.
This resulted in a single ngbPropBurn variable being added to each
model.

The quality of the final model was assessed via inspection of residuals
and a 10-fold cross-validation. In each fold, we used 90\% of the data
to train the model and 10\% to test the model; we calculated the Cox
Snell R2, the global deviance (i.e.~twice minus log-likelihood; {[}CITE:
Stasinopoulos-2017{]}) and root mean squared error, and compared their
values against those from the full model. Finally, we analysed the
partial effect of each covariate and covariate interaction on four
components of severity: the expected mean severity (Equation~\ref{eq-6})
and variance (Appendix 1), the estimated probability of observing
stand-replacement (p1; Equation~\ref{eq-5}) and the estimated
probability of observing no severity (p0; Equation~\ref{eq-4}). In the
case of interactions between two continuous variables (e.g.~Broadleaf *
Height), we fixed one of them (Height) at the 5\%, 50\% and 95\%
quantile values to show effects of the first covariate (Broadleaf) at
average and extreme values of the second covariate. We also show maps of
expected mean severity and probability of stand-replacement for the two
fires that showed the highest and the lowest agreement between observed
and predicted severity (measured as the Pearson correlation value, r).

\section{Results}\label{results}

The best fitting and most parsimonious ZOIB regression model had
different sets of covariates for each parameter of the ZOIB distribution
and the random effect of FireID on the intercepts of all parameter
models:

\begin{equation}\phantomsection\label{eq-7}{
\begin{aligned}
\text{Severity} &\sim \text{ZOIB}(\mu, \sigma, \nu, \tau) \\
\mu &\sim \text{Broadleaf} * \text{Height} + \text{Lari} + \text{CTI} + \text{SMRwet} + \text{ngbPropBurn960m} + (1 \mid \text{FireID}) \\
\sigma &\sim \text{Broadleaf} * \text{Understorey} + \text{SMRwet} + \text{ngbPropBurn1920m} + (1 \mid \text{FireID}) \\
\nu &\sim \text{FlamConif} + \text{SRR} + \text{SMRwet} + \text{ngbPropBurn30m} + (1 \mid \text{FireID}) \\
\tau &\sim \text{Broadleaf} * \text{Height} + \text{Broadleaf} * \text{Understorey} + \text{Lari} * \text{Height} + \text{Lari} * \text{Understorey} \\
     &\quad + \text{isForest} + \text{CTI} + \text{SMRwet} + \text{SRR} + \text{ngbPropBurn30m} + (1 \mid \text{FireID})
\end{aligned}
}\end{equation}

where ``*'' denotes the additive effects and interaction of any two
covariates. Despite our attempts to summarise fire weather in different
ways, no fire weather covariates had significant effects on fire
severity and were kept in the final model. Model residuals confirmed
that the ZOIB distribution adequately fit the data and that the
parameter models were correctly specified (Fig. S3.1 Appendix 3). The
overall fit of the model was very satisfactory, considering the large
variation and poor covariate signal observed during exploratory
analysis. The model explained ca. 46.3\% of the total variance, as
measured by the Cox and Snell R2, and had a global deviance of 483,785
(red dots in Fig. 2); across the 10 folds used in cross-validation, R2
values varied less than 0.1\% and global deviance varied between 59,000
and 60,750. The model performed well at high severity values, but
overestimated severity at lower levels (Fig. S3.2 Appendix 3). We also
tested the model's capacity to accurately predict the observed severity
classes (see Appendix 2 for how continuous severity predictions were
classified). In this case, the model accuracy ranged from
51.82\%-86.67\% across the six classes, being highest for classes 0 and
5 (86.67\% and 68.98\%, respectively); overall model accuracy was
58.14\%. The model tended to overpredict stand-replacement in
intermediate classes (Table 6). However, if we consider the
disproportionate number of stand-replacement observations in the
dataset, the model performed relatively well in predicting classes 1 to
4, with classes 2 and 3 being the least accurately predicted (Table 6).

Since each component of severity (i.e.~mean, variance, p0 and p1) is
linked to more than one parameter, each is affected by the group of
covariates entering the models of the parameters they are linked to.
Below, we focus on these effects in forested cells, except when the
variable coding forested and non-forested cells (isForested) is a
covariate in the model. Responses in non-forested cells were, in
general, qualitatively similar but quantitatively dampened, which could
be due to the much smaller sample size (425,120 forested cells
vs.~12,437 non-forested cells; Figs. S3.2-S2.6 Appendix 3).

\subsection{Expected mean severity}\label{expected-mean-severity}

Expected mean severity was most affected by the proportion of bordering
neighbours burnt (ngbPropBurn30m; Tables 2 and 3), to which it showed a
sigmoidal positive response (Fig. 3). At intermediate scales
(ngbPropBurn960m -- third strongest effect), the proportion of burnt
neighbour cells had a negative effect on mean severity, but it is worth
noting that at this scale there were always \textgreater60\% burnt
neighbour cells (Fig. 3). Stand height had the second strongest effect
on mean severity and a significant interaction with Broadleaf cover
(Broadleaf) and tamarack cover (Lari; Fig. 3, Tables 2 and 3). The
taller and more Broadleaf-dominated the stands, the lower was the
expected mean severity especially in the presence of understorey
vegetation, which also interacted with Broadleaf cover (Fig. 3). On the
other hand, higher values of tamarack dominance increased mean severity
in tall stands, but decreased it in short stands (Fig. 3). Presence of
forest (isForest) and terrain roughness (SRR; fourth and fourteenth
strongest effects) all led to higher mean severity values, as did the
presence of understory when considered independently of Broadleaf cover.
Surprisingly, flammable conifer (FlamConif) cover only had a small
positive effect on mean severity (twelfth strongest effect; Tables 2 and
3, Fig. 3). Moisture indicators (topographic soil wetness, CTI, and soil
moisture derived from forest inventory plots, SMR; the eighth and tenth
most important effects) had negative effects on mean severity (Tables 2
and 3, Fig. 3).

\subsection{Probabilities of stand-replacement and no
severity}\label{probabilities-of-stand-replacement-and-no-severity}

The probability of stand-replacement (p1) was affected by the same
covariates as mean severity (with the exception of burnt neighbours at
intermediate scale, ngbPropBurn960m), to which it showed qualitatively
similar responses (Table 2, Fig. 4). For instance, higher proportions of
bordering neighbours burnt (ngbPropBurn30m), tall stands with less
Broadleaf and more tamarack cover all increased p1, while higher
moisture values (CTI and SMR) decreased p1 values (Fig. 4). As with mean
severity, flammable conifer cover had a small positive effect on p1. We
note, however, that the effect of the proportion of burnt bordering
neighbours did not increase p1 above 0.75.

The probability of observing no severity, i.e.~no post-fire mortality
(p0), was affected by the same covariates as p1, but with qualitatively
opposite responses (Table 2, Fig. 5). The exception to this pattern was
terrain roughness, which had a positive effect on both p0 and p1 (Table
2, Fig. 5). The effect of flammable conifer cover was slightly more
pronounced, and negative, on p0 and the response of p0 to the proportion
of bordering neighbours burnt (ngbPropBurn30m) covered the full range of
p0 values, from 0 to 1 (Fig. 5).

\subsection{Variance of severity}\label{variance-of-severity}

As with the other components of severity, variance values were most
affected by the proportion of bordering neighbours burnt
(ngbPropBurn30m), but the response was hump-shaped (Tables 2 and 3, Fig.
6). Variance approached zero at very low proportions of burnt neighbours
and was maximised around 55\% of burnt neighbours. As the proportion of
burnt neighbours increased, variance decreased again but did not
approach zero. Severity variance was also affected by the proportion of
burnt neighbours calculated at intermediate scale (ngbPropBurn960m) and
at large scale (ngbPropBurn1920m). In these two cases the response was
linear, but positive at intermediate scale and slightly negative at
large scale (Fig. 6). Height, in interaction with Broadleaf and tamarack
cover, also had a strong effect on variance, with taller, more
Broadleaf-dominated stands and less tamarack-dominated stands showing
larger variance in fire severity. The effects of Broadleaf and tamarack
cover on variance were also modulated by the presence/absence of
understorey vegetation. In the case of tamarack cover this interaction
was weak, but in Broadleaf-dominated stands variance was larger when
understorey vegetation was present (Tables 2 and 3, Fig. 7). Otherwise,
in general, the presence of understorey vegetation contributed to
decrease variance in severity. Flammable conifer cover also decreased
variance, as did the presence of forest (Fig. 7) and terrain roughness
(Fig. 6). On the other hand, higher moisture values were related with
higher severity variability (CTI, Fig. 6, and SMR, Fig. 7).

\subsection{Severity maps}\label{severity-maps}

Two Alberta fires, Buffalo Head and Livock, had the highest agreement
values between observed and expected mean severity (r = 0.83 and r =
0.76, respectively). The Buffalo Head fire occurred in the Boreal Forest
region and the Livock fire in the Foothills. A visual inspection of
observed and predicted severity values for these two fires shows a very
good match between observed and predicted unburnt patches (Fig. 8). In
these patches, the median expected mean severity was below 0.05
(approximately 0.03 and 0.04 for Buffalo Head and Livock, respectively)
and the median probability of a stand-replacing fire was below 0.01
(approximately 0.006 and 0.009). Stand-replacement predictions were less
accurate. In patches where observations indicated complete mortality,
the median expected mean severity was 0.92 (Buffalo Head) and 0.86
(Livock) and the median probability of a stand-replacing fire was 0.71
(Buffalo Head) and 0.57 (Livock). Mapping estimated variance showed
higher variability in severity in the interface between unburnt and
burnt patches (Fig. 8). Similar patterns were observed in fires where
there was lower agreement between observations and predictions, with
several unburnt patches being well captured by the model but very high
severity patches not appearing distinct from other more intermediate
values (Fig. S3.8); yet, in these fires, high values of variance were
observed across most of the burnt area and not just at the interface
between burnt and unburnt patches (Fig. S3.8).

\section{Discussion}\label{discussion}

We present a statistical model of fire severity that overcomes the
challenges associated with modelling zero-and-one-inflated proportional
data while accounting, at least partially, for distinct ways in which
different biophysical drivers affect fire severity. Our
zero-and-one-inflated beta (ZOIB) regression model based on a GAMLSS
framework partitions the effects of distinct sets of covariates on
average severity, on the probability of stand-replacement or
non-burning, and on the variability of severity, all of which are
important from a forest and fire management perspective. For instance,
the probabilities of a stand-replacement event (or, conversely, of an
unburnt patch) can be combined with estimates of average fire severity
to weight predictions in areas where fire damage is expected to be
highest (or lowest). Similarly, modelling variance as a function of
potentially distinct covariates not only improves model fit, but also
provides information on the processes that drive variance {[}CITE:
Goldstein-2005{]} and that in this case may decrease certainty in
estimated fire severity. This can inform managers and decision makers
about areas where management actions may need to remain flexible or
``prepare for the worst'', but also allow them to reduce this variance
by controlling factors that drive it. In our case, higher variance was
in general associated with the same covariate values that led to lower
average severity and to more unburnt patches (Figs. 3, 5-7), and
appeared highest around predicted unburnt patches (Fig. 8). Forest and
fire management may want to focus on these interface areas when
assessing opportunities to manipulate fire attributes to lower the risk
of high severity fires, e.g.~by increasing the cover of mature broadleaf
stands and controlling the abundance and composition of ladder fuels
(the understorey).

Fire severity is driven by bottom-up (fuel and topography) and top-down
drivers (climate; {[}REF{]}). Weather, in particular, has been shown to
override the effects of fuel when conditions are extremely warm and dry
(often called extreme fire weather conditions; {[}REF{]}). Yet, despite
our efforts, weather variables did not show any significant effects on
severity and were not retained in the final model. There might be
several explanations as to why this happened. First, the a-spatial
nature of our weather data may be behind our inability to detect
significant weather effects on severity. In our model, fire ID was added
as a random effect and the lack of a signal of weather on severity
suggests that most of the weather variability is likely tied to fire ID.
It is possible that in our dataset, weather is playing a more important
role in explaining differences between fires and fire properties like
size, the proportion of patches of a given severity and the level of
heterogeneity in severity within a fire perimeter (things we did not
analyse). Second, temporal variation in weather likely influences fire
severity as the fire progresses, but we did not have information on fire
progression that could have allowed temporally pairing fire weather and
severity. These data limitations may also explain why our model
underpredicted stand-replacement, which is probably largely driven by
extreme fire weather conditions {[}REF{]}. Finally, it is also possible
that, in our case, within-fire spatial severity patterns were indeed
more influenced by topographic and fuel factors than by variation in
weather conditions during fire. Ferster et al.'s {[}REF{]} analysis
using multinomial logistic models also showed fuel characteristics,
followed by topography and soil moisture, were the most important
drivers of severity patterns within fire, while fire weather (via
drought code, DC) was a driver of severity variability between fires.
Pre-fire vegetation structure and composition have also been shown to be
important drivers of fire severity in mixed-conifer forests in Montana
({[}CITE: Belote-2015{]}) and Sierra Nevada ({[}CITE: Lydersen-2016{]}),
USA, where pre-fre species composition, size and tree density all
affected post-fire mortality levels.

{[}CITE: Belote-2015{]} also show a strong effect of pre-fire community
composition, tree characteristics (species, i.e.~traits, and size) on
post-fire tree mortality at the community level. {[}CITE:
Lydersen-2016{]} show strong effects of tree density on post-fire
mortality. They noted that plots with greeter shrub cover tended to have
higher severity class. With respect to pre-fire forest characteristics,
severity classes varied more in terms of pre-fire tree density than in
species composition - they only saw that hardwood basal area was greater
in moderate and high severity plots, compared to low severity plots.

had relatively good performance in predicting both continuous and
class-based values of severity, when compared with previous similar
work. For instance, {[}CITE: Fang-2018{]} reported overall accuracy
levels 52.4-67.1\% when using boosted regression trees to predict fire
severity (in three classes) in 2000 and 2010 fires in Chinese boreal
forests, while {[}CITE: Ferster-2016{]} reported an 39\% overall
accuracy level when using multinomial logistic models to predict fire
severity in six classes using the same dataset as we used here. Unlike
what we expected and what was found by {[}CITE: Ferster-2016{]}, fire
weather did not have a significant effect on driving maximum severity
values or any other component of severity (mean, variance, or no
severity), despite our attempts to summarise and include fire weather
variables in different ways.

\begin{itemize}
\tightlist
\item
  weather variables not retained

  \begin{itemize}
  \tightlist
  \item
    {[}CITE: Belote-2015{]} show a strong effect of pre-fire community
    composition, tree characteristics (species, i.e.~traits, and size)
    on post-fire tree mortality at the community level.
  \item
    {[}CITE: Lydersen-2016{]} show strong effects of tree density on
    post-fire mortality. They noted that plots with greeter shrub cover
    tended to have higher severity class. With respect to pre-fire
    forest characteristics, severity classes varied more in terms of
    pre-fire tree density than in species composition - they only saw
    that hardwood basal area was greater in moderate and high severity
    plots, compared to low severity plots.
  \item
    It may also be that our weather data was limited because it was not
    spatially explicit and because we don't have information about the
    temporal development of fires. Using a soil moisture variable like
    Fang et al did could be a potential solution. This could explain why
    our model was not as good in detecting very high severity (notice
    that at very high severity values the model underpredicted severity)
  \end{itemize}
\item
  Our model reveals strong effects of fuels on fire severity which, in a
  way, is good news because we can control fuels but not the weather.

  \begin{itemize}
  \tightlist
  \item
    presence of understorey - wasn't always significant but throughout
    it seems to have a general effect of decreasing severity and
    increasing its variance.
  \item
    the effect of stand height, on the other hand depended on species
    composition and vertical structure

    \begin{itemize}
    \tightlist
    \item
      lower severity if tall stands were dominated by Broadleaf and
      there was an understorey
    \item
      higher severity if tall stands were dominated by tamarack
    \end{itemize}
  \end{itemize}
\item
  this model is useful in a simulation context and more transferable

  \begin{itemize}
  \tightlist
  \item
    At the local scale (disregarding differences between distinct fire
    events and looking at cells within a perimeter) landscape variables
    may be playing a more important role in explaining differences in
    severity levels.
  \item
    we can use it to predict severity while taking into account the
    effect of drivers on extreme events (zero and ones)
  \end{itemize}
\item
  In all cases, including our own, intermediate severity levels were the
  hardest to predict with severity often being overestimated (Table 4
  and Fig. S3.2; {[}CITE: Fang-2018{]}, Fig. 4; {[}CITE:
  Ferster-2016{]}, tbl. 6).
\item
  different spatial scales act on different components. potential
  artifact of very large buffer size and being close to the fire
  perimeter
\end{itemize}

\section*{References}\label{references}
\addcontentsline{toc}{section}{References}

\section{Tables}\label{tables}

Table 1. Full set of covariates initially investigated as potential
drivers of severity. Covariates in bold entered were selected for model
fitting, but not all were kept in the final model. All vegetation
covariates, except for `isForest', `SMR', `SMRwet' and `understorey'
refer to the overstory layer.

\begin{longtable}[]{@{}
  >{\raggedright\arraybackslash}p{(\linewidth - 6\tabcolsep) * \real{0.2500}}
  >{\raggedright\arraybackslash}p{(\linewidth - 6\tabcolsep) * \real{0.2500}}
  >{\raggedright\arraybackslash}p{(\linewidth - 6\tabcolsep) * \real{0.2500}}
  >{\raggedright\arraybackslash}p{(\linewidth - 6\tabcolsep) * \real{0.2500}}@{}}
\caption{}\label{tbl-1}\tabularnewline
\toprule\noalign{}
\begin{minipage}[b]{\linewidth}\raggedright
Type
\end{minipage} & \begin{minipage}[b]{\linewidth}\raggedright
Name
\end{minipage} & \begin{minipage}[b]{\linewidth}\raggedright
Description
\end{minipage} & \begin{minipage}[b]{\linewidth}\raggedright
Unit
\end{minipage} \\
\midrule\noalign{}
\endfirsthead
\toprule\noalign{}
\begin{minipage}[b]{\linewidth}\raggedright
Type
\end{minipage} & \begin{minipage}[b]{\linewidth}\raggedright
Name
\end{minipage} & \begin{minipage}[b]{\linewidth}\raggedright
Description
\end{minipage} & \begin{minipage}[b]{\linewidth}\raggedright
Unit
\end{minipage} \\
\midrule\noalign{}
\endhead
\bottomrule\noalign{}
\endlastfoot
Vegetation & \textbf{Abies} & Percentage of a species or generic group
of species that contributes to the species composition of a cell. Must
add up to 100\% across species in a stand. \emph{Abies balsamea} --
balsam fir; \emph{Abies lasiocarpa} -- subalpine fir & \% \\
Vegetation & \textbf{Lari} & \emph{Larix laricina} -- tamarack & \% \\
Vegetation & \textbf{FlamConif} & \emph{Picea glauca} -- white spruce;
\emph{Picea engelmannii} -- Engelmann's spruce; \emph{Picea mariana} --
black spruce; \emph{Pinus banksiana} -- jack pine; \emph{Pinus contorta}
-- lodgepole pine & \% \\
Vegetation & \textbf{Broadleaf} & \emph{Betula papyrifera} -- paper
birch; \emph{Populus balsamifera} -- balsam poplar; \emph{Populus
tremuloides} -- trembling aspen & \% \\
Vegetation & Crown closure (lower/upper) & Crown closure (lower and
upper bounds) & \% \\
Vegetation & \textbf{isForest} & Is the cell forested? This includes
forested wetlands & TRUE/FALSE \\
Vegetation & \textbf{Height} & Stand height class & meters \\
Vegetation & Non-forest type & Type of land-cover in non-forested cells
& `Infrastructure', `grassland', `industry', `river', `low shrub', `tall
shrub' \\
Vegetation & Origin & Average initiation year of dominant and codominant
species in the cell, determined to the nearest year or decade. &
years \\
Vegetation & \textbf{SMR} & Soil moisture regime converted to an ordered
variable (`dry' = 1, `mesic' = 2, `moist' = 3, `wet' = 4, `aquatic' = 5)
& unitless \\
Vegetation & \textbf{SMRwet} & Soil moisture regime converted to a
binary variable indicating if the soil is wet or aquatic (values 1-3 =
FALSE, values 4 and 5 = TRUE) & TRUE/FALSE \\
Vegetation & \textbf{Understorey} & Presence/absence of understorey &
TRUE/FALSE \\
Weather & {[}mean/min/max/cv/range{]}Temp & Mean, minimum, maximum,
coefficient of variation and range of temperature summarised across all
fire weather days, only the first three days, or only the first day (for
which CV = 0) & degrees Celsius \\
Weather & {[}mean/min/max/cv/range{]}RH & Mean, minimum, maximum,
coefficient of variation and range of relative humidity summarised
across all fire weather days, only the first three days, or only the
first day (for which CV = 0) & \% \\
Weather & {[}mean/min/max/cv/range{]}WS & Mean, minimum, maximum,
coefficient of variation and range of wind speed summarised across all
fire weather days, only the first three days, or only the first day (for
which CV = 0) & km/h \\
Weather & {[}mean/min/max/cv/range{]}Rain & Mean, minimum, maximum,
coefficient of variation and range of rain precipitation summarised
across all fire weather days, only the first three days, or only the
first day (for which CV = 0) & mm \\
Weather & {[}mean/min/max/cv/range{]}FFMC & Mean, minimum, maximum,
coefficient of variation and range of fine fuel moisture code summarised
across all fire weather days, only the first three days, or only the
first day (for which CV = 0) & unitless \\
Weather & {[}mean/min/max/cv/range{]}DMC & Mean, minimum, maximum,
coefficient of variation and range of duff moisture code summarised
across all fire weather days, only the first three days, or only the
first day (for which CV = 0) & unitless \\
Weather & {[}mean/min/max/cv/range{]}DC & Mean, minimum, maximum,
coefficient of variation and range of drought code summarised across all
fire weather days, only the first three days, or only the first day (for
which CV = 0) & unitless \\
Weather & {[}mean/min/max/cv/range{]}ISI & Mean, minimum, maximum,
coefficient of variation and range of initial spread index summarised
across all fire weather days, only the first three days, or only the
first day (for which CV = 0) & unitless \\
Weather & {[}mean/min/max/cv/range{]}BUI & Mean, minimum, maximum,
coefficient of variation and range of build up index summarised across
all fire weather days, only the first three days, or only the first day
(for which CV = 0) & unitless \\
Weather & {[}mean/min/max/cv/range{]}FWI & Mean, minimum, maximum,
coefficient of variation and range of fire weather index summarised
across all fire weather days, only the first three days, or only the
first day (for which CV = 0) & unitless \\
Topography & \textbf{Elevation} & Ground elevation in m above sea level
& m \\
Topography & \textbf{Slope} & Ground slope in percent & \% \\
Topography & \textbf{Aspect} & Ground aspect in degrees & degrees \\
Topography & \textbf{SAR} & Surface Area and Ratio - measure of terrain
roughness (area/cos(slope in radians)) & unitless \\
Topography & \textbf{SRR} & Surface relief ratio - measure of terrain
roughness & unitless \\
Topography & \textbf{CTI} & Compound topographic index - measure of soil
wetness & unitless \\
Other & \textbf{ngbPropBurns{[}30/60/120/240/480/960/1920{]}} &
Proportion of burnt neighbours for each cell in the analysis at
different buffer sizes (30m, 60m, 120m, 240m, 480m, 960m, 1920m) &
proportion \\
Other & \textbf{FireID} & Fire name & -- \\
\end{longtable}

Table 2. Final zero-and-one-inflated beta regression model. A different
set of covariates was fit for the parameters \(\mu\), \(\sigma\),
\(\nu\) and \(\tau\) of the zero-and-one-inflated beta distribution
(ZOIB(\(\mu\), \(\sigma\), \(\nu\), \(\tau\))). Covariates were
re-scaled between 0 and 1 to compare estimated effects. t-values in bold
are significant at p-value \textless{} 0.05.

\begin{longtable}[]{@{}lllll@{}}
\caption{}\label{tbl-2}\tabularnewline
\toprule\noalign{}
Parameter & Covariates & Estimate & Std. Error & t-value \\
\midrule\noalign{}
\endfirsthead
\toprule\noalign{}
Parameter & Covariates & Estimate & Std. Error & t-value \\
\midrule\noalign{}
\endhead
\bottomrule\noalign{}
\endlastfoot
\(\mu\) & (Intercept) & 0.18 & 0.004 & \textbf{44.29} \\
& ngbPropBurn960m & -0.972 & 0.003 & \textbf{-38.8} \\
& Lari & 0.068 & 0.003 & \textbf{25.47} \\
& Broadleaf & -0.105 & 0.004 & \textbf{-24.71} \\
& CTI & -0.063 & 0.003 & \textbf{-22.22} \\
& SMRwet & -0.09 & 0.008 & \textbf{-10.86} \\
& Height & 0.028 & 0.004 & \textbf{8.02} \\
& Broadleaf:Height & 0.013 & 0.002 & \textbf{5.14} \\
\(\sigma\) & (Intercept) & -0.23 & 0.004 & \textbf{-54.77} \\
& ngbPropBurn1920m & 0.114 & 0.003 & \textbf{43.84} \\
& Understorey (TRUE) & 0.062 & 0.005 & \textbf{11.95} \\
& SMRwet & 0.042 & 0.006 & \textbf{7.48} \\
& Broadleaf:Understorey (TRUE) & 0.029 & 0.005 & \textbf{5.66} \\
& Broadleaf & 0.02 & 0.004 & \textbf{5.45} \\
\(\nu\) & (Intercept) & -2.83 & 0.015 & \textbf{182.44} \\
& ngbPropBurn30m & -2.126 & 0.01 & \textbf{-210.73} \\
& SMRwet & 0.262 & 0.02 & \textbf{13.36} \\
& FlamConif & -0.069 & 0.009 & \textbf{-7.78} \\
& SRR & 0.042 & 0.009 & \textbf{4.77} \\
\(\tau\) & (Intercept) & -0.059 & 0.027 & \textbf{-2.17} \\
& ngbPropBurn30m & 0.583 & 0.008 & \textbf{76.34} \\
& Height & -0.265 & 0.006 & \textbf{-40.88} \\
& isForest & 0.787 & 0.027 & \textbf{29.00} \\
& Understorey (TRUE) & 0.178 & 0.008 & \textbf{23.31} \\
& Broadleaf:Height & -0.084 & 0.004 & \textbf{-20.01} \\
& CTI & -0.059 & 0.004 & \textbf{-13.73} \\
& Broadleaf:Understorey (TRUE) & -0.096 & 0.008 & \textbf{-12.63} \\
& Height:Lari & 0.067 & 0.009 & \textbf{7.42} \\
& SMRwet & -0.079 & 0.012 & \textbf{-6.77} \\
& SRR & 0.023 & 0.004 & \textbf{5.58} \\
& Lari:Understorey (TRUE) & 0.039 & 0.008 & \textbf{4.70} \\
& Broadleaf & -0.026 & 0.006 & \textbf{-4.33} \\
& Lari & -0.009 & 0.014 & -0.65 \\
\end{longtable}

Table 3. Importance of model covariate towards the variation of each
component of severity. Covariates entering the models of any group of
parameters affecting a given component of severity are shown in
decreasing order of their absolute t-value (interpreted as covariate
importance; see Table 2 for values). We noted in parentheses when a
covariate affects a severity component via two or more parameters.

\textbf{Expected mean, E(Y\^{}r) -- ZOIB parameters: \(\mu\), \(\nu\),
\(\tau\)}

Covariate importance (decreasing): ngbPropBurn30m (via \(\nu\) and
\(\tau\)); Height (via \(\mu\) and \(\tau\)); ngbPropBurn960m; isForest;
Lari (via \(\mu\) and \(\tau\)); Broadleaf (via \(\mu\) and \(\tau\));
Understorey; CTI (via \(\mu\) and \(\tau\)); Broadleaf:Height (via
\(\mu\) and \(\tau\)); SMRwet (via \(\mu\), \(\nu\) and \(\tau\));
Broadleaf:Understorey; FlamConif; Lari:Height; SRR (via \(\nu\) and
\(\tau\)); Lari:Understorey.

\textbf{p0 and p1, P(Y=0) and P(Y=1) -- ZOIB parameters: \(\nu\),
\(\tau\)}

Covariate importance (decreasing): ngbPropBurn30m (via \(\nu\) and
\(\tau\)); Height; isForest; Understorey; Broadleaf:Height; CTI; SMRwet
(via \(\nu\) and \(\tau\)); Broadleaf:Understorey; FlamConif;
Lari:Height; SRR (via \(\nu\) and \(\tau\)); Lari:Understorey;
Broadleaf; Lari.

\textbf{Expected variance, Var(Y) -- ZOIB parameters: \(\mu\),
\(\sigma\), \(\nu\), \(\tau\)}

Covariate importance (decreasing): ngbPropBurn30m (via \(\nu\) and
\(\tau\)); ngbPropBurn1920m; Height (via \(\mu\) and \(\tau\));
ngbPropBurn960m; isForest; Lari (via \(\mu\) and \(\tau\)); Broadleaf
(via \(\mu\), \(\sigma\) and \(\tau\)); Understorey (via \(\sigma\) and
\(\tau\)); CTI (via \(\mu\) and \(\tau\)); Broadleaf:Height (via \(\mu\)
and \(\tau\)); SMRwet (via \(\mu\), \(\sigma\), \(\nu\) and \(\tau\));
Broadleaf:Understorey (via \(\sigma\) and \(\tau\)); FlamConif;
Lari:Height; SRR (via \(\nu\) and \(\tau\)); Lari:Understorey.

Table 4. Confusion matrix for the final ZOIB regression model. True
positives (diagonal) are underlined and the observed percent occurrence
of each class is shown in parentheses.

\begin{longtable}[]{@{}
  >{\raggedright\arraybackslash}p{(\linewidth - 12\tabcolsep) * \real{0.1429}}
  >{\raggedright\arraybackslash}p{(\linewidth - 12\tabcolsep) * \real{0.1429}}
  >{\raggedright\arraybackslash}p{(\linewidth - 12\tabcolsep) * \real{0.1429}}
  >{\raggedright\arraybackslash}p{(\linewidth - 12\tabcolsep) * \real{0.1429}}
  >{\raggedright\arraybackslash}p{(\linewidth - 12\tabcolsep) * \real{0.1429}}
  >{\raggedright\arraybackslash}p{(\linewidth - 12\tabcolsep) * \real{0.1429}}
  >{\raggedright\arraybackslash}p{(\linewidth - 12\tabcolsep) * \real{0.1429}}@{}}
\toprule\noalign{}
\begin{minipage}[b]{\linewidth}\raggedright
Predicted severity class
\end{minipage} & \begin{minipage}[b]{\linewidth}\raggedright
0 (9\%)
\end{minipage} & \begin{minipage}[b]{\linewidth}\raggedright
1 (6\%)
\end{minipage} & \begin{minipage}[b]{\linewidth}\raggedright
2 (6\%)
\end{minipage} & \begin{minipage}[b]{\linewidth}\raggedright
3 (7\%)
\end{minipage} & \begin{minipage}[b]{\linewidth}\raggedright
4 (9\%)
\end{minipage} & \begin{minipage}[b]{\linewidth}\raggedright
5 (62\%)
\end{minipage} \\
\midrule\noalign{}
\endhead
\bottomrule\noalign{}
\endlastfoot
0 & 76\% & 4\% & 5\% & 5\% & 4\% & 2\% \\
1 & 11\% & 18\% & 13\% & 10\% & 8\% & 3\% \\
2 & 4\% & 14\% & 13\% & 12\% & 7\% & 5\% \\
3 & 2\% & 12\% & 15\% & 13\% & 7\% & 5\% \\
4 & 3\% & 12\% & 14\% & 14\% & 13\% & 9\% \\
5 & 5\% & 41\% & 40\% & 47\% & 62\% & 76\% \\
\end{longtable}

\begin{description}
\item[Observed severity class (\% occurrence) is shown across columns.]
\{\#tbl-4\}
\end{description}

\section{Figures}\label{figures}

\textbf{Figure 1.} The 36 fire perimeters (in black) with severity,
topography, pre-fire vegetation and fire weather data included in the
zero-one-inflated beta regression models. Canadian Ecozones (top legend)
and Natural Regions of Alberta (bottom legend and colour-outlined
polygons) are also shown.

\emph{{[}FIGURE PLACEHOLDER: fig-1{]}}

\textbf{Figure 2.} Cross-validation results and model fit. The Cox and
Snell R2 and global deviance were calculated for each fold of a k-fold
cross-validation (k = 10) and their distributions are shown in the
boxplots. The red dots show the R2 and global deviance of the model
fitted with the full dataset.

\emph{{[}FIGURE PLACEHOLDER: fig-2{]}}

\textbf{Figure 3.} Covariate partial effects on mean severity in
forested cells. The expected mean of a ZOIB regression is affected by
parameters \(\mu\), \(\nu\) and \(\tau\); hence, the figure shows the
partial effects of covariates entering the \(\mu\), \(\nu\) and \(\tau\)
parameter models (see Tables 2 and 3), where covariate values were
back-transformed to their original scale. The plots show marginal mean
predictions (i.e.~excluding the random effect of FireID).

\emph{{[}FIGURE PLACEHOLDER: fig-3{]}}

\textbf{Figure 4.} Covariate partial effects on expected probabilities
of observing maximum severity, p1 (i.e.~100\% mortality) in forested
cells. Given that p1 in a ZOIB regression is affected by parameters
\(\nu\) and \(\tau\), the figure shows the partial effects of all
covariates in the \(\nu\) and \(\tau\) parameter models (see Tables 2
and 3), where covariate values were back-transformed to their original
scale. The plots show marginal mean predictions (i.e.~excluding the
random effect of FireID).

\emph{{[}FIGURE PLACEHOLDER: fig-4{]}}

\textbf{Figure 5.} Covariate partial effects on expected probabilities
of observing no severity, p0 (i.e.~a zero - 0\% mortality) in forested
cells. Given that p0 in a ZOIB regression is affected by parameters
\(\nu\) and \(\tau\), the figure shows the partial effects of all
covariates in the \(\nu\) and \(\tau\) parameter models (see Tables 2
and 3), where covariate values were back-transformed to their original
scale. The plots show marginal mean predictions (i.e.~excluding the
random effect of FireID).

\emph{{[}FIGURE PLACEHOLDER: fig-5{]}}

\textbf{Figure 6.} Covariate partial effects on the expected variance of
severity in forested cells. Given that the expected variance is affected
by all parameters of the ZOIB regression, the figure shows the partial
effects of all covariates included in the ZOIB regression model (see
Tables 2 and 3), where covariate values were back-transformed to their
original scale. The plots show marginal mean predictions (i.e.~excluding
the random effect of FireID). See Figure 7 for remaining covariates.

\emph{{[}FIGURE PLACEHOLDER: fig-6{]}}

\textbf{Figure 7.} Continued from Figure 6.

\emph{{[}FIGURE PLACEHOLDER: fig-7{]}}

\textbf{Figure 8.} Maps of observed and predicted severity for the a)
Buffalo Head and b) Livock fires, for which there was high agreement
between predictions and observations. Predicted severity maps show
expected mean severity (second column) and the probability of
stand-replacement (i.e., maximum severity, p1; third column).

\emph{{[}FIGURE PLACEHOLDER: fig-8{]}}

\section{Appendix 1}\label{appendix-1}

The variance of a ZOIB can be calculated by expanding {[}CITE:
Ospina-2010{]}'s formulation using the ZOIB(alpha, gamma, \(\mu\),
\(\sigma\)) to the ZOIB(\(\mu\), \(\sigma\), \(\nu\), \(\tau\))
parameterisation implemented in \texttt{gamlss} {[}CITE:
Stasinopoulos-2007{]}; see also {[}CITE: Rigby-2019{]}), as follows:

\[
\text{Var}(Y) = \alpha V_1 + (1 - \alpha) V_2 + \alpha (1 - \alpha)(\gamma - \mu)^2
\]

where,

\[
\gamma = \frac{\tau (1 + \nu + \tau)}{\nu + \tau} \cdot (1 + \nu + \tau)
\]

\[
\alpha = \frac{\nu + \tau}{1 + \nu + \tau}
\]

\[
V_1 = \gamma (1 - \gamma)
\]

\[
V_2 = \frac{\mu (1 - \mu)}{1 + \sigma}
\]

Note that \(V_2\) is the variance of the Beta process (same as
Var(\(\mu\)) in Equation~\ref{eq-1} in the main text).

\section{Appendix 2}\label{appendix-2}

\subsection{Supplementary methods}\label{supplementary-methods}

\subsubsection{Exploratory Principal Components Analyses and Spearman
Rank
correlations}\label{exploratory-principal-components-analyses-and-spearman-rank-correlations}

The following Principal Components Analysis (PCAs) were calculated using
vegetation and topography covariates used to characterise stands. On
each PCA we overlaid severity and fire weather variables using the
\texttt{env.fit} function of the \texttt{vegan} R package.
\texttt{env.fit} tries to maximise the correlation value between a focus
variable vector (SEVERITY) and the PCA axes, allowing to visualise
correlations between the variables entering the PCA and the variable(s)
added post-hoc.

\emph{{[}FIGURE PLACEHOLDER: fig-s2-1{]}}

\textbf{Figure S2.1.} PCA of stand characteristics, coloured by fire and
with overlaid fire weather variables and severity. Right side plot is a
zoom-in on the variable vectors. Covariates of stand characteristics
(full arrows) include species groups' cover (\emph{Abies} spp, `Abie',
\emph{Larix laricina}, `Lari', flammable conifers, `FlamConif', and
broadleaf species, `Broadleaf'), stand height, crown closure, presence
of understory, soil moisture regime, elevation, slope and aspect. Fire
weather and severity (SEV\_CONT) variables were added a posteriori using
\texttt{env.fit} (dashed arrows). Only the mean and coefficient of
variation (prefixed with `cv') of each fire weather variable was added
to the PCA. Please see Table 1 in main text for all variable names and
description.

\emph{{[}FIGURE PLACEHOLDER: fig-s2-2{]}}

\textbf{Figure S2.2.} PCA of stand characteristics, coloured by soil
moisture regime (SMR) and with overlaid fire weather variables and
severity. Right side plot is a zoom-in on the variable vectors. See the
legend of Figure S2.1 for further details and Table 1 in main text for
all variable names and description.

\subsubsection{Conversion of expected mean severity to severity
classes}\label{conversion-of-expected-mean-severity-to-severity-classes}

A confusion matrix was calculated to more directly compare our results
to Ferster et al.'s {[}REF{]} (Table 6). This required comparing the
predictions of the final ZOIB regression model (fitted to the full
dataset using scaled covariates) with observed severity classes. For
this, we converted the expected mean severity values (calculated
following Equation~\ref{eq-6} in the main text) to ordinal classes,
whose intervals were defined using the quantile values (\(Q_i\)) of the
cumulative probabilities, \(P_i\), of each severity class j, calculated
as:

\[
P_i = \sum_{j=1}^{i} x_j, \quad \text{for } i \in \{1, 2, \ldots, N\}
\]

where \(x_j\) is the observed probability of each class (calculated as
the number of cells with a given severity class divided by the total
sample size) and N is the total number of classes (6). The resulting
quantile values were approximately
\(Q_{1,\ldots,6} = \{0.09, 0.16, 0.22, 0.29, 0.38, 1.00\}\) and the
predicted class intervals in \% mortality were: class 0 {[}0, 46\%);
class 1 {[}46\%, 65\%); class 2 {[}65\%, 72\%); class 3 {[}72\%, 82\%);
class 4 {[}82\%, 99\%); and class 5 {[}99\%, 100{]}. These class
intervals were then used to convert the expected mean severity values to
predicted severity classes. The confusion matrix in Table 4 (main text)
is a cross-tabulation of the observed severity classes (see Fire
severity data section in Materials and Methods) and predicted severity
classes per cell, obtained with \texttt{confusionMatrix} function in the
\texttt{caret} R package {[}CITE: Kuhn-2022{]}.

\section{Appendix 3}\label{appendix-3}

\subsection{Supplementary tables and
figures}\label{supplementary-tables-and-figures}

\emph{{[}FIGURE PLACEHOLDER: fig-s3-1{]}}

\textbf{Figure S3.1.} Residual inspection plots for the final ZOIB
regression model following {[}CITE: Stasinopoulos-2017{]}.

\emph{{[}FIGURE PLACEHOLDER: fig-s3-2{]}}

\textbf{Figure S3.2.} Expected mean severity vs.~observed severity. The
blue line shows the predicted linear regression fit and the dashed line
represents a 1:1 line. The Pearson correlation value is annotated in the
upper left corner.

\emph{{[}FIGURE PLACEHOLDER: fig-s3-3{]}}

\textbf{Figure S3.3.} Covariate partial effects on expected mean
severity in non-forested cells. Given that the expected mean value of a
ZOIB regression is affected by parameters \(\mu\), \(\nu\) and \(\tau\),
the figure shows the partial effects of all covariates in the \(\mu\),
\(\nu\) and \(\tau\) parameter models (see Tables 2 and 3), where
covariate values were back-transformed to their original scale. The
plots show marginal mean predictions (i.e.~excluding the random effect
of FireID).

\emph{{[}FIGURE PLACEHOLDER: fig-s3-4{]}}

\textbf{Figure S3.4.} Covariate partial effects on expected
probabilities of observing maximum severity, p1 (i.e.~one - 100\%
mortality) in non-forested cells. Given that p1 in a ZOIB regression is
affected by parameters \(\nu\) and \(\tau\), the figure shows the
partial effects of all covariates in the \(\nu\) and \(\tau\) parameter
models (see Tables 2 and 3), where covariate values were
back-transformed to their original scale. The plots show marginal mean
predictions (i.e.~excluding the random effect of FireID).

\emph{{[}FIGURE PLACEHOLDER: fig-s3-5{]}}

\textbf{Figure S3.5.} Covariate partial effects on expected
probabilities of observing no severity, p0 (i.e.~a zero - 0\% mortality)
in non-forested cells. Given that p0 in a ZOIB regression is affected by
parameters \(\nu\) and \(\tau\), the figure shows the partial effects of
all covariates in the \(\nu\) and \(\tau\) parameter models (see Tables
2 and 3), where covariate values were back-transformed to their original
scale. The plots show marginal mean predictions (i.e.~excluding the
random effect of FireID).

\emph{{[}FIGURE PLACEHOLDER: fig-s3-6{]}}

\textbf{Figure S3.6.} Covariate partial effects on expected severity
variance in non-forested cells. Given that the expected variance is
affected by all parameters in the ZOIB regression, the figure shows the
partial effects of all covariates included in the ZOIB regression model
(see Tables 2 and 3), where covariate values were back-transformed to
their original scale. The plots show marginal mean predictions
(i.e.~excluding the random effect of FireID). See Figure S3.7 for
remaining covariates.

\emph{{[}FIGURE PLACEHOLDER: fig-s3-7{]}}

\textbf{Figure S3.7.} Continued from Figure S3.6.

\emph{{[}FIGURE PLACEHOLDER: fig-s3-8{]}}

\textbf{Figure S3.8.} Maps of observed and predicted severity for the a)
Contest \#2 and b) Woodman fires, for which there was the lowest
agreement between observed and predicted values. Predicted severity maps
show expected mean severity (second column) and the predicted
probability of observing maximum severity (p1, third column).


\bibliography{references.bib}



\end{document}
